
Soit @@M0@@~; on a :
@@M1@@
@@M2@@
l'égalité du centre résultant de IV 2.12.2.

Ce résultat permet de prouver 1.2 par récurrence noethérienne sur l'ensemble des sous-schémas fermés de @@M3@@. Si on suppose que pour tout sous-schéma fermé @@M4@@ de @@M5@@, distinct de @@M6@@, @@M7@@ est surjective, alors le raisonnement précédent montre qu'on peut se ramener à prouver 1.2 pour un schéma de base irréductible et affine, quitte à remplacer @@M8@@ par un ouvert affine irréductible. De plus, le diagramme
@@M9@@
où @@M10@@ est surjectif d'après 1.1, montre qu'on peut également supposer @@M11@@ réduit.

On se place donc sur un schéma de base intègre et affine. On pose @@M12@@, et on note @@M13@@ le corps des fractions de @@M14@@; si @@M15@@ est le point générique de @@M16@@, alors @@M17@@.

Soit @@M18@@. Par hypothèse, @@M19@@ est engendré par les images des @@M20@@; si @@M21@@ est l'image de @@M22@@ dans @@M23@@, on peut trouver des entiers @@M24@@ tels que :
@@M25@@
